\documentclass[UTF8,zihao=-4,oneside]{ctexrep}

% 论文主线：单目 USB 相机采集目标信息，PC 通过 USB/TTL 直接控制六个舵机。
% 版式和学校模板单独处理，本文件先保证论文内容、实验状态和图片素材对齐。
\usepackage[a4paper,left=2.8cm,right=2.6cm,top=2.8cm,bottom=2.6cm]{geometry}
\usepackage{amsmath}
\usepackage{booktabs}
\usepackage{float}
\usepackage{graphicx}
\usepackage{hyperref}
\usepackage{tabularx}

\graphicspath{{../../prepared_images/}}

\hypersetup{
  colorlinks=true,
  linkcolor=black,
  urlcolor=black,
  citecolor=black
}

\newcommand{\paperfigure}[4]{
  \begin{figure}[H]
    \centering
    \includegraphics[width=#4\textwidth,keepaspectratio]{#3}
    \caption{#1}
    \label{#2}
  \end{figure}
}

\newcommand{\placeholderfigure}[4]{
  \begin{figure}[H]
    \centering
    \fbox{
      \begin{minipage}[c][#4][c]{0.86\textwidth}
        \centering
        #3
      \end{minipage}
    }
    \caption{#1}
    \label{#2}
  \end{figure}
}

\title{\heiti 海参采捕场景下双关节六绳驱动软体机器人视觉伺服控制研究}
\author{梁淑轩}
\date{\today}

\begin{document}

\maketitle

\begin{abstract}
针对海参采捕作业中人工潜水风险高、传统刚性执行机构适应复杂海底环境能力不足以及软体目标易受损等问题，本文围绕双关节六绳驱动软体机器人的视觉伺服控制开展研究。系统采用单目 USB 相机获取目标图像，PC 侧完成图像采集、目标检测、误差计算和控制指令生成，并通过 USB/TTL 串口链路直接驱动舵机牵引拉索，使软体关节产生弯曲响应。

在机构设计方面，本文参考模块化绳驱连续体关节的设计思路，对原始四绳单关节方案进行改型，形成“单关节三单元、三根绳驱动，整机双关节、六根绳驱动”的实验平台。功能分工上，第一关节承担单目视觉伺服对准任务，第二关节以手动控制方式承担吸取范围扩展和姿态补偿作用。样机验证以第二关节为先，第二关节有效长度为 426 mm，三根驱动绳孔位半径为 40 mm，完整装配第一关节后扩展为双关节六绳系统。

围绕上述硬件主线，本文完成了单目相机棋盘格标定、海参数据集构建、YOLO26n 检测模型训练、视频目标跟踪、图像中心误差提取以及舵机串口直控基础测试。标定阶段共使用 51 幅有效棋盘格图像，鱼眼模型平均重投影误差为 0.0569 px；海参检测数据集包含训练集 48 幅、验证集 12 幅，最佳验证结果达到 mAP@0.5 为 0.993，mAP@0.5:0.95 为 0.742。实验结果说明，视觉感知链路能够为图像误差到舵机控制量的闭环映射提供输入。

本文围绕机构设计、视觉感知、数据训练、运动学建模和舵机执行链路开展研究，建立从图像平面误差到拉索控制量的实现路径，为双关节六绳驱动软体机器人视觉伺服实验提供结构基础和控制方法。

\noindent\textbf{关键词：}海参采捕；绳驱软体机器人；视觉伺服；YOLO 目标检测；USB/TTL 直控
\end{abstract}

\tableofcontents

\chapter{绪论}

\section{研究背景及意义}

海参具有较高的经济价值，但当前采捕过程仍然高度依赖人工潜水作业。人工方式不仅劳动强度大、成本高，而且在复杂海底环境下存在明显安全风险。对于面向规模化作业的海洋装备而言，如何在尽量不破坏海参本体的前提下实现稳定对准与柔顺采捕，是一项同时具有工程价值和应用价值的问题。

海参通常附着或停留于泥沙、礁石等复杂底质环境中，周围障碍物多，目标本体柔软且易受刺激。传统刚性机械臂与夹爪在此类场景下面临两个典型困难：一是构型刚性强，在狭窄环境中的贴近性差；二是接触过程缺少柔顺性，容易造成目标滑脱或损伤。绳驱连续体与软体机器人具备柔性变形能力和较好的环境适应性，为海参采捕类近场作业提供了新的执行机构思路。

在软体机器人具备柔顺性的基础上，仅依靠开环动作仍难以应对目标位置不确定、视角变化和环境扰动等问题，因此有必要引入视觉伺服控制。对本课题而言，单目相机可以提供目标中心、边界框和图像平面误差等关键信息；若进一步建立“图像误差—拉索控制量—关节变形—目标再观测”的闭环，就能够形成一条适合现有实验条件、可逐步验证的技术路径。

\section{国内外相关研究现状}

\subsection{水下刚性机械臂研究现状}

传统水下机械臂多采用刚性连杆和关节构成，驱动形式以液压、电机或电液混合为主，优点是负载能力强、控制模型成熟，已经在载人潜器、ROV 和海上作业平台中得到应用。但对于海参采捕这类近底、狭窄、软体目标接触任务，刚性机械臂需要较高的定位精度和接触控制能力，末端夹持或碰撞容易对目标造成损伤。

\paperfigure{典型水下刚性机械臂结构图（引自冯家琪，2023）}{fig:rigid-manipulator-ref}{fig1_1_rigid_manipulators.jpg}{0.82}

刚性水下机械臂通常追求较大的负载能力和较高的末端定位精度，适合水下打捞、取样、阀门操作和设备维护等任务。现有系统多将电机、液压缸或传动机构布置在关节附近，整机刚度较高，但密封、重量和结构体积会随自由度数量增加而上升。对于海底养殖环境中的海参采捕任务，目标附近常存在泥沙、礁石和养殖设施，刚性机构在贴近目标时容易受到空间约束，末端执行器与海参接触时也需要额外的柔顺控制。

因此，刚性机械臂虽然为水下作业提供了成熟的工程基础，但其结构特点与海参柔性目标的近距离采捕需求并不完全匹配。本文将刚性机械臂作为对比背景，重点研究更适合软体目标接触和狭窄环境贴近的绳驱软体执行机构。

\subsection{绳驱连续体与软体机器人研究现状}

连续体机器人通过柔性杆、弹性骨架、弹簧或模块化关节形成连续弯曲运动，能够在没有传统刚性转动副的情况下改变末端位置和姿态。绳驱连续体机器人利用多根拉索牵引柔性本体产生弯曲，驱动电机可集中布置在基座处，有利于减小末端质量，并能提高水下结构的密封便利性。该类机器人已在医疗内窥操作、管道检测、航空航天维护和水下仿生操作等场景中得到研究。

国外研究中，绳驱连续体结构常采用超弹性材料、弹簧骨架或柔性支撑杆作为中心骨架，并沿圆周方向布置三根或四根驱动绳，通过绳长差形成全方位弯曲运动。部分方案能够实现较大弯曲角和多段串联运动，但也存在驱动箱体积较大、绳路摩擦复杂、控制参数多和水下密封困难等问题。

国内相关研究多围绕仿象鼻、仿章鱼触手和柔性连续机械臂展开，重点包括结构设计、正逆运动学、工作空间分析和样机实验。冯家琪关于水下仿象鼻绳驱连续型机械臂的研究将机械臂运动学划分为驱动空间、关节空间和操作空间，分别讨论绳长变化量、关节变量和末端位姿之间的映射关系。该组织方式适合绳驱连续体机构的建模，也为本文三绳单关节运动学分析提供了章节组织参考。

北京航空航天大学翟士民等设计的球铰嵌套式连续型机械臂采用多个球铰单元串联，中间设置弹性橡胶垫圈，并由每节三根绳索牵引形成弯曲运动。该类方案属于典型的欠驱动连续型机械臂，特点是结构简单、可靠性较高，并且能够用较少驱动绳实现单关节两个弯曲自由度。该思路与本文采用三绳单关节的设计取向一致，说明三绳方案在连续体机构中具有可借鉴的工程基础。

模块化连续体关节设计能够兼顾结构加工便利性和柔顺性调节能力，对本科阶段实体样机搭建具有直接参考价值。Dewi 等（2024）提出的模块化绳驱连续体机器人方案为本文机构构型提供了概念来源。本文并不直接照搬该文的四绳单关节结构，而是在现有舵机数量、加工条件和调试难度约束下进行三绳化和双关节化改型。

\paperfigure{模块化绳驱连续体机器人结构图（引自 Dewi 等，2024）}{fig:continuum-robot-ref}{fig1_2_tdcr_reference.png}{0.88}

从现有研究看，多绳驱动可以提高关节姿态调节能力，但同时会增加绳路耦合、驱动协调和控制建模难度。本文结合海参采捕任务的柔顺接触需求和现有加工条件，采用单关节三单元、三根绳驱动的结构形式，并进一步串联形成双关节六绳驱动软体机器人。

\subsection{目标检测研究现状}

水下视觉感知受到光照衰减、悬浮颗粒散射、蓝绿偏色和背景纹理干扰影响，传统阈值分割或人工特征方法在复杂底质环境下稳定性不足。深度学习目标检测方法可以直接从图像中学习目标形状、纹理和尺度特征，使海参目标的框选、中心定位和连续跟踪具备工程实现基础。

目标检测的输出通常包括类别、置信度和边界框位置。对本文实验平台而言，检测框中心点可以直接转化为图像平面坐标，检测框宽高可以反映目标尺度变化，置信度可以用于判断当前检测结果是否可靠。相比只给出分类结果的图像识别方法，目标检测更适合作为视觉伺服控制的前端输入。

海参目标形态柔软、边界不规则，且实验图像中可能出现姿态变化、局部遮挡和背景纹理干扰。因此，检测模型不仅需要识别目标类别，还需要稳定输出目标中心位置。本文采用 YOLO26n 完成单类别海参检测训练，其目的不是单独构建一个目标检测演示，而是为控制器提供可连续读取的目标中心和尺度信息。

\subsection{视觉伺服研究现状}

视觉伺服方法通常分为基于图像特征的 IBVS、基于位姿的 PBVS 及其混合策略。结合本课题硬件条件与代码基础，本文以检测框中心、图像中心和误差滤波为核心变量，建立从图像误差到舵机增量控制的简化闭环，并在此基础上引入绳驱连续体运动学建模和协调控制。

IBVS 直接利用图像平面特征构造误差，控制目标是使图像特征逐步收敛到期望位置；PBVS 需要估计目标与相机之间的相对位姿，再在三维空间中计算控制量。对于本文的海参采捕实验平台，检测框中心与图像中心之间的像素差可以直接形成误差变量，适合先建立图像平面内的对准控制。第二关节承担吸取范围扩展作用时，可与第一关节对准动作配合，使视觉误差修正与末端作用范围调整分开处理。

\subsection{本文切入点}

已有研究通常分别强调绳驱连续体机构设计、轨迹控制或视觉伺服算法，而围绕“单目视觉 + 双关节六绳驱动软体机器人 + 海参采捕场景”这一组合的工程研究相对较少。本文面向本科毕业设计实验条件，采用单目 USB 相机、PC 侧视觉处理和 USB/TTL 舵机直控方案，研究从海参图像检测、目标中心误差提取到绳驱软体关节控制量生成的实现方法。

\section{主要研究内容}

本文主要围绕以下四项内容展开：
\begin{enumerate}
  \item 完成双关节六绳驱动软体机器人实验平台总体方案设计，明确两关节的功能分工、机构构型、驱动方式和控制链路。
  \item 完成单目相机标定、去畸变、海参数据采集、标注和 YOLO26n 检测模型训练，为视觉伺服提供稳定视觉输入。
  \item 结合串口驱动脚本与舵机执行结构，提出从图像误差到绳驱控制量的简化映射思路，并以三绳单关节模型作为运动学验证对象。
  \item 通过标定结果、检测结果、目标跟踪预览和舵机直控能力，对实验平台视觉输入与执行链路进行验证分析。
\end{enumerate}

\chapter{双关节六绳驱动软体机器人系统方案与结构设计}

\section{系统总体方案}

系统采用“单目视觉感知 + PC 侧集中控制 + USB/TTL 串口直控六舵机 + 双关节六绳驱动软体机器人执行”的总体方案。相机负责采集海参目标图像，PC 完成检测推理、目标中心提取与误差计算，再通过 USB/TTL 转换链路向舵机下发控制命令，驱动拉索完成关节弯曲与姿态调整。两个关节在功能上进行分工：第一关节用于接收单目视觉误差并完成图像平面对准，第二关节用于扩大吸取范围和进行末端姿态补偿；样机验证中第二关节先采用手动方式完成方向和幅值验证。

\paperfigure{双关节六绳驱动软体机器人总体结构设计图}{fig:overall-structure}{fig2_1.png}{0.92}

该方案的特点是硬件路径短、调试对象明确，能够减少中间环节对实验判断的干扰。本文以 PC 直控舵机作为主要实现路径，便于在实验平台搭建过程中直接观察图像误差、舵机角度和关节运动之间的对应关系。

\section{双关节六绳驱动机构设计}

本文机构参考 Dewi 等提出的模块化连续体关节思路，但没有直接沿用其原始结构。原始方案中一个关节通常由多个模块单元和四根拉索共同驱动；结合样机加工条件、舵机数量和调试复杂度，本文将单关节调整为三个结构单元、三根绳驱动，并将整机设计为两个关节串联，因此完整执行端共使用六根绳和六个舵机。

单关节采用三根绳驱动的原因在于该方案结构简单且有效。三根驱动绳在圆周方向相隔 \(120^\circ\) 布置时，已经能够使单关节在任意方向产生弯曲，满足平面误差修正和末端范围调整所需的两个弯曲自由度。与四绳方案相比，三绳方案减少了舵机数量、走线数量和绳路摩擦点，也降低了装配调试、方向标定和控制分配矩阵整定的复杂度。对当前本科阶段实验平台而言，三绳单关节能够在控制难度和运动能力之间取得较稳妥的平衡。

\paperfigure{软体机器人单元结构实物图}{fig:unit-photo}{fig2_2.jpg}{0.65}

每个单元近似为圆柱体，高度为 142 mm，该高度包含上下两个圆盘，每个圆盘厚度为 3 mm。每个关节由 3 个单元组成，因此单个关节总高度为 426 mm；两个关节全部安装后的设计总高度为 852 mm。第一关节三根驱动绳孔位半径为 20 mm，第二关节三根驱动绳孔位半径为 40 mm。这样的尺寸分配使第一关节更适合承担视觉伺服中的方向修正任务，第二关节则可在手动控制下提供更明显的末端偏移范围。

从俯视方向定义，第一关节驱动点为点一、点二、点三，相邻夹角为 \(120^\circ\)，对应绳一、绳二、绳三。第二关节驱动点整体相对第一关节顺时针偏置 \(40^\circ\)，对应点四、点五、点六，也就是绳四、绳五、绳六。样机已安装第二关节及其对应三根驱动绳，第二关节限位位置为 426 mm；完整安装第一关节时，第二关节限位位置为满足连接结构调整到 429 mm。

\paperfigure{六绳孔位俯视布置示意图}{fig:rope-layout}{fig2_4_rope_layout.png}{0.78}

\paperfigure{六舵机驱动部分实物图}{fig:servo-drive-photo}{fig2_3.jpg}{0.72}

驱动部分采用舵机直接牵引拉索的形式布置于机器人基部，便于集中供电、统一固定和串口调试。该布置适合桌面验证和参数调整，也便于从单关节验证扩展到完整双关节联动。

\section{硬件组成与控制链路}

\begin{table}[H]
  \centering
  \caption{实验平台硬件组成}
  \label{tab:hardware-platform}
  \begin{tabularx}{\textwidth}{>{\raggedright\arraybackslash}p{0.24\textwidth} >{\raggedright\arraybackslash}p{0.30\textwidth} X}
    \toprule
    模块 & 配置 & 在本实验平台中的作用 \\
    \midrule
    视觉输入 & 单目 USB 相机 & 获取海参目标图像，输出检测框中心和图像误差 \\
    标定对象 & 10\(\times\)7 内角点棋盘格，方格边长 20 mm & 求取相机内参与畸变参数，支撑去畸变处理 \\
    执行机构 & 双关节六绳驱动软体机器人 & 第一关节用于视觉伺服对准，第二关节用于手动扩展吸取范围；实验先验证第二关节运动学和控制接口 \\
    舵机链路 & PC 经 USB/TTL 串口直接控制舵机 & 下发位置控制命令并读取角度、电压、电流和功率状态 \\
    控制软件 & Python 脚本 & 完成检测推理、误差提取、运动学映射和串口指令发送 \\
    \bottomrule
  \end{tabularx}
\end{table}

本文执行端选用 Fashion Star HA8-U25-M UART 总线舵机\footnote{\url{https://fashionstar.com.hk/wiki/servo/uart/datasheet/pdf/Fashion-Star-HA8-U25-M.pdf}}。该舵机工作电压为 6.0--8.4 V，采用 UART/TTL 半双工通信，支持 9600 bps 到 1 Mbps 波特率、0--254 ID 寻址、单圈角度模式、多圈角度模式和阻尼模式；在 7.4 V 条件下，其静态堵转扭矩为 2.45 N\(\cdot\)m，额定扭矩为 0.44 N\(\cdot\)m。对本文六绳软体机器人而言，其可取之处主要体现在总线布线简洁、六个舵机可统一寻址、可读取位置和运行状态、多圈角度模式适合绕线轮连续收放，以及温度、电压、堵转、功率和电流保护有利于降低拉索卡滞时的过载风险。

\section{本章小结}

本章说明了双关节六绳驱动软体机器人的总体方案、机构改型依据、两关节职能和硬件链路。系统采用“单目相机 + PC 集中控制 + 舵机直接牵引拉索”的实验路线，为图像误差到多舵机协调动作的闭环控制提供硬件基础。

\chapter{单目视觉检测与目标误差提取方法}

\section{单目相机标定与去畸变}

为提高图像误差计算的稳定性，本文首先完成了单目 USB 相机的棋盘格标定与去畸变处理。所用 USB 相机原始画面存在明显广角畸变，图像边缘位置会出现几何拉伸和直线弯曲现象。如果直接使用畸变图像计算目标中心偏差，同一实际位移在图像中心和边缘会对应不同像素变化量，进而影响视觉伺服控制中的误差判断。因此，去畸变不是单纯的图像美化步骤，而是建立稳定图像误差坐标系的前提。

\paperfigure{USB 相机规格与畸变参数截图}{fig:camera-spec}{fig3_0_camera_spec.jpg}{0.70}

相机成像可近似看作三维点经透视投影后落到图像平面。设相机坐标系下空间点为 \(P_c=(X_c,Y_c,Z_c)^T\)，理想归一化成像坐标为
\begin{equation}
  x=\frac{X_c}{Z_c},\quad y=\frac{Y_c}{Z_c}
\end{equation}
若不考虑镜头畸变，像素坐标由相机内参矩阵 \(K\) 决定：
\begin{equation}
  \begin{bmatrix}
    u\\v\\1
  \end{bmatrix}
  =
  \begin{bmatrix}
    f_x & 0 & c_x\\
    0 & f_y & c_y\\
    0 & 0 & 1
  \end{bmatrix}
  \begin{bmatrix}
    x\\y\\1
  \end{bmatrix}
\end{equation}
实际广角镜头会使归一化坐标发生径向畸变。以鱼眼模型为例，令 \(r=\sqrt{x^2+y^2}\)，\(\theta=\arctan r\)，畸变后的角度可表示为
\begin{equation}
  \theta_d=\theta(1+k_1\theta^2+k_2\theta^4+k_3\theta^6+k_4\theta^8)
\end{equation}
畸变坐标由 \(\theta_d\) 反推到图像平面：
\begin{equation}
  x_d=\frac{\theta_d}{r}x,\quad y_d=\frac{\theta_d}{r}y
\end{equation}
相机标定的目的就是求出 \(f_x,f_y,c_x,c_y\) 以及畸变系数 \(k_1,k_2,k_3,k_4\)。去畸变时，根据这些参数把畸变图像中的像素点反向映射到理想成像位置，使图像直线关系和目标中心位置尽量恢复到统一几何坐标系中。对视觉伺服而言，这一步直接影响目标中心误差的可信度。

标定流程采用 10\(\times\)7 内角点棋盘格，方格边长为 20 mm。程序首先在多幅不同姿态、不同位置的棋盘格图像中检测角点，并将角点坐标细化到亚像素级；随后建立棋盘格平面点在标定板坐标系中的三维坐标，结合图像角点坐标求解相机内参矩阵和鱼眼畸变系数。实际运行时不逐帧重新求解畸变参数，而是根据标定结果预先生成映射表，再利用重映射完成实时正畸，从而兼顾校正效果和视频处理速度。

\paperfigure{棋盘格标定图像采集过程}{fig:checkerboard-capture}{fig3_1.jpg}{0.82}

\paperfigure{标定阶段的棋盘格角点提取结果}{fig:checkerboard-corners}{fig3_2.jpg}{0.82}

\begin{table}[H]
  \centering
  \caption{单目相机标定结果}
  \label{tab:calibration-result}
  \begin{tabular}{ll}
    \toprule
    参数 & 数值 \\
    \midrule
    有效标定图像数量 & 51 \\
    标定模型 & fisheye \\
    图像分辨率 & 1920\(\times\)1080 \\
    \(f_x / f_y\) & 1064.830 / 1060.592 \\
    \(c_x / c_y\) & 902.598 / 527.784 \\
    RMS & 0.9420 \\
    平均重投影误差 / px & 0.0569 \\
    \bottomrule
  \end{tabular}
\end{table}

\paperfigure{相机原始画面示意图}{fig:raw-camera-view}{fig3_3.jpg}{0.82}

\paperfigure{去畸变预览结果示意图}{fig:undistort-preview}{fig3_4.png}{0.82}

从标定结果可以看出，相机模型已经得到可用于图像校正的参数。平均像素误差较小，说明角点检测与参数求解结果具有较好一致性。图 \ref{fig:raw-camera-view} 和图 \ref{fig:undistort-preview} 对比显示，正畸后画面边缘弯曲得到明显修正，图像中心和边缘区域的几何关系更加统一。由于本文关注图像平面内的对准关系，去畸变处理主要用于减少边缘区域几何误差对目标中心计算的影响。

\section{数据集构建与海参检测模型训练}

围绕海参目标检测任务，本文完成了数据采集、标注、训练集划分和模型训练等工作。目标检测模块采用 YOLO26n 单阶段检测模型，其基本思想是在一次前向推理中同时预测目标边界框、类别概率和置信度。与传统先生成候选区域再分类的方法相比，单阶段检测模型更适合视频流推理和实时误差提取。

数据集包含一个目标类别 \texttt{seacucumber}，训练集 48 幅、验证集 12 幅。该数据规模能够支撑模型训练、推理测试和视觉误差提取。对本文而言，模型训练本身也是实验平台建设的一部分，其作用是将原始图像输入转化为可被控制器使用的目标中心坐标、检测框尺度和置信度。

\paperfigure{海参检测训练数据采集示意图}{fig:data-collection}{fig3_5.jpg}{0.70}

\paperfigure{海参检测数据标注结果示意图}{fig:data-labeling}{fig3_6.png}{0.82}

标注阶段使用矩形框标出图像中的海参目标，并保存为 YOLO TXT 格式。每个标注文件记录目标类别编号以及归一化后的中心点坐标、框宽和框高。视觉伺服控制直接使用检测框中心，因此标注质量不仅影响检测精度，也会影响目标中心位置的稳定性。标注过程以单类别海参为主，数据扩充时仍需保持同一类别命名和标注规范，避免训练集前后版本不一致。

\begin{table}[H]
  \centering
  \caption{海参检测数据集划分情况}
  \label{tab:dataset-split}
  \begin{tabularx}{\textwidth}{>{\raggedright\arraybackslash}p{0.22\textwidth} >{\raggedright\arraybackslash}p{0.18\textwidth} >{\raggedright\arraybackslash}p{0.18\textwidth} X}
    \toprule
    数据划分 & 图像数量 & 标注框数量 & 说明 \\
    \midrule
    训练集 & 48 & 155 & 用于模型训练与参数更新 \\
    验证集 & 12 & 38 & 用于精度评估与早停判断 \\
    总计 & 60 & 193 & 单类别海参检测数据集 \\
    \bottomrule
  \end{tabularx}
\end{table}

表 \ref{tab:dataset-split} 为数据集统计结果。随着实验平台继续采集复杂背景、不同光照和不同目标姿态样本，表中图像数量和标注框数量需要同步更新。

训练过程中使用 YOLO26n 预训练权重，输入尺寸设置为 832，批大小为 8，训练上限为 150 轮，并采用 patience=30 的早停策略。实际训练在第 88 轮结束，说明模型在本数据规模下已进入相对稳定的收敛阶段。

\paperfigure{海参检测训练过程截图}{fig:training-process}{fig3_7.png}{0.82}

\paperfigure{训练过程中损失函数与检测指标变化曲线}{fig:training-curves}{extra_fig_results.png}{0.88}

\begin{table}[H]
  \centering
  \caption{海参检测训练与验证结果}
  \label{tab:detection-metrics}
  \begin{tabularx}{\textwidth}{>{\raggedright\arraybackslash}p{0.30\textwidth} >{\raggedright\arraybackslash}p{0.18\textwidth} >{\raggedright\arraybackslash}p{0.18\textwidth} X}
    \toprule
    指标 & 数值 & 对应轮次 & 说明 \\
    \midrule
    最佳 mAP@0.5 & 0.993 & 57 & 本小样本验证集上的最高检测准确率 \\
    最佳 mAP@0.5:0.95 & 0.742 & 58 & 更严格 IoU 条件下的最好结果 \\
    最后一轮 Precision & 0.878 & 88 & 训练结束时模型仍保持可用检测能力 \\
    最后一轮 Recall & 0.943 & 88 & 对目标召回相对充分，适合继续作为控制输入基础 \\
    \bottomrule
  \end{tabularx}
\end{table}

由于验证集规模为 12 幅图像、38 个标注框，上述结果属于本数据版本下的模型训练结论。对视觉伺服任务而言，这些指标的意义主要在于说明模型能够稳定输出目标中心，为误差计算提供基础输入。

图 \ref{fig:training-curves} 展示了训练过程中损失函数和检测指标的变化。训练前期损失下降较快，说明模型能够从本数据集中学习到海参目标的主要纹理和形状特征；中后期 mAP 曲线趋于稳定，说明在本小样本数据集上模型已经基本收敛。图 \ref{fig:label-distribution} 反映了标注框中心位置和尺度分布，可用于判断样本是否过度集中在某一画面区域。图 \ref{fig:val-prediction} 给出了验证批次预测结果，重点观察检测框是否覆盖目标主体、是否存在漏检和误检，而不仅是查看最终数值指标。

\paperfigure{数据集标注框中心与尺度分布图}{fig:label-distribution}{extra_fig_labels.jpg}{0.70}

\paperfigure{验证批次预测结果可视化}{fig:val-prediction}{extra_fig_valpred.jpg}{0.82}

\section{目标跟踪与图像误差提取}

在推理程序实现上，脚本在检测结果基础上引入基于 \texttt{track} 接口的目标跟踪功能，并在显示界面中叠加图像中心十字线、目标框中心点和目标到图像中心的偏差向量。目标跟踪的作用是在连续帧之间保持目标编号和检测框关联，减少单帧检测漏检、框位置突变和目标切换对控制输入的影响。

具体流程为：摄像头或视频帧首先经过可选去畸变处理，再送入 YOLO 模型进行检测与跟踪；程序根据检测框坐标计算目标中心点，并与图像中心点作差，得到水平和竖直方向的像素误差。显示界面中的红色连线表示目标中心到图像中心的偏差向量，数值形式的 \(dx,dy\) 则作为控制器输入。这样可以将视觉输出转化为可供控制器使用的误差变量，而不是停留在“只画检测框”的展示层面。

\paperfigure{海参目标检测效果预览图}{fig:detection-preview}{fig3_8.jpg}{0.72}

\paperfigure{目标跟踪与图像误差输出预览图}{fig:tracking-error-preview}{fig3_9.png}{0.86}

设输入图像分辨率为 \(W \times H\)，图像中心定义为
\begin{equation}
  u_0 = \frac{W}{2}, \quad v_0 = \frac{H}{2}
\end{equation}
对于检测或跟踪得到的目标框，其中心坐标记为 \((u, v)\)，则图像平面误差定义为
\begin{equation}
  e_x = u - u_0, \quad e_y = v_0 - v
\end{equation}
其中，\(e_x\) 表示目标相对图像中心的水平偏差，\(e_y\) 表示目标相对图像中心的竖直偏差。程序采用“图像坐标向右为正、向上为正”的误差约定，便于把误差与关节弯曲方向联系起来。

\section{本章小结}

本章围绕视觉伺服所需的图像输入，完成了 USB 相机畸变分析、单目相机标定、实时去畸变、海参检测模型训练、标注流程、跟踪界面构建及图像误差定义。上述内容将原始图像转化为目标中心和像素误差，为控制器计算舵机控制量提供视觉变量。

\chapter{面向六舵机绳驱系统的视觉伺服控制方法}

\section{控制目标与实现约束}

视觉伺服控制的目标是在海参目标被稳定检测到的条件下，根据图像平面误差 \(e_x, e_y\) 计算舵机增量控制量，使绳驱软体机器人调整关节弯曲方向，并逐步将目标移动到图像中心附近。本文将控制变量首先定义在图像平面内，再通过舵机与拉索关系转化为软体关节的弯曲运动。

该控制问题受到三类约束。第一，执行机构为绳驱软体机器人，关节变形由多根拉索共同决定，不能直接套用刚性机械臂的 D-H 建模和关节独立控制方法。第二，检测与跟踪输出存在位置抖动和目标切换，若直接驱动舵机，容易造成拉索频繁往复。第三，系统采用 PC 直控舵机路径，需要在串口时序、命令限幅和异常处理上保持简洁可调。

\paperfigure{视觉伺服控制系统总体方案图}{fig:visual-servo-control-scheme}{fig4_1_control_system_scheme.pdf}{0.92}

\section{单关节运动学模型与基本假设}

绳驱连续体机器人的运动学分析通常围绕驱动空间、关节空间和操作空间三者之间的映射关系展开。驱动空间对应舵机角度和绳长变化量，关节空间对应弯曲角和弯曲方向角，操作空间对应末端位置和姿态。本文将任意一个三绳驱动关节看作单关节模型进行推导，第一关节和第二关节均可采用同一模型，只需代入各自的结构参数和绳孔角位置。

单关节模型建立在以下假设基础上：关节中心线近似满足常曲率弯曲；柔性支撑段轴向伸长相对较小，关节弧长近似保持为 \(L\)；三根驱动绳沿圆周方向布置并随关节一起弯曲；驱动绳孔位半径 \(R\) 在弯曲过程中保持不变。上述假设使绳长变化量、弯曲角和末端位姿之间可以用几何投影关系表达。

建立单关节基坐标系 \(B=\{O,x,y,z\}\)，其中 \(O\) 位于该关节近端基座圆盘中心，\(+z\) 轴沿关节未弯曲时的中心轴线，\(+x\) 轴选取为第一根驱动绳所在的参考方向，\(+y\) 轴按右手定则确定。单关节只定义一个近端基座圆盘和一个远端末端圆盘；完整双关节串联时，第一关节直接连接在第二关节末端圆盘处。

\paperfigure{单关节运动学坐标系建系图}{fig:single-joint-coordinate}{fig4_2_single_joint_coordinate.png}{0.78}

为避免与实物编号混淆，本节用 \(i=1,2,3\) 表示单个关节内的三根驱动绳。对于第一关节，三根绳对应实物绳一、绳二、绳三；对于第二关节，三根绳对应实物绳四、绳五、绳六。记三根驱动绳在本关节坐标系下的角位置为 \(\alpha_i\)，相邻驱动绳间隔为 \(120^\circ\)。单关节弯曲角记为 \(\theta\)，弯曲方向角记为 \(\phi\)，第 \(i\) 根绳的长度变化量记为 \(\Delta l_i\)，舵机绕线轮角度记为 \(q_i\)。符号约定为 \(\Delta l_i>0\) 表示收紧或缩短，\(q_i>0\) 表示舵机按收紧方向旋转。

\section{驱动空间与关节空间正逆运动学}

驱动空间到绳长变化量的关系由绕线轮半径决定：
\begin{equation}
  \Delta l_i = r_w q_i,\quad i\in\{1,2,3\}
\end{equation}

关节空间到驱动空间的正运动学为：已知 \(\theta,\phi\)，求三根绳长变化量和舵机角度。关节向方向 \(\phi\) 弯曲时，半径为 \(R\)、角位置为 \(\alpha_i\) 的驱动绳长度变化量等于弯曲量 \(R\theta\) 在该驱动绳方向上的投影，因此
\begin{equation}
  \Delta l_i = R\theta\cos(\phi-\alpha_i)
\end{equation}
三根驱动绳可写成一组方程：
\begin{equation}
  \left\{
  \begin{aligned}
    \Delta l_1 &= R\theta\cos(\phi-\alpha_1)\\
    \Delta l_2 &= R\theta\cos(\phi-\alpha_2)\\
    \Delta l_3 &= R\theta\cos(\phi-\alpha_3)
  \end{aligned}
  \right.
\end{equation}
再由
\begin{equation}
  q_i=\frac{\Delta l_i}{r_w}
\end{equation}
得到对应舵机角度。

驱动空间到关节空间的逆运动学为：已知 \(\Delta l_1,\Delta l_2,\Delta l_3\)，求 \(\theta,\phi\)。先定义
\begin{equation}
  k_x=R\theta\cos\phi,\quad k_y=R\theta\sin\phi
\end{equation}
则每根驱动绳长度变化量可写为
\begin{equation}
  \Delta l_i=k_x\cos\alpha_i+k_y\sin\alpha_i
\end{equation}
三根绳组成矩阵形式：
\begin{equation}
  \begin{bmatrix}
    \Delta l_1\\
    \Delta l_2\\
    \Delta l_3
  \end{bmatrix}
  =
  \begin{bmatrix}
    \cos\alpha_1 & \sin\alpha_1\\
    \cos\alpha_2 & \sin\alpha_2\\
    \cos\alpha_3 & \sin\alpha_3
  \end{bmatrix}
  \begin{bmatrix}
    k_x\\
    k_y
  \end{bmatrix}
\end{equation}
利用三根绳在圆周方向等角布置的性质，可以得到闭式解
\begin{equation}
  k_x = \frac{2}{3}\sum_{i=1}^{3}\Delta l_i\cos\alpha_i,\quad
  k_y = \frac{2}{3}\sum_{i=1}^{3}\Delta l_i\sin\alpha_i
\end{equation}
于是关节变量为
\begin{equation}
  \theta=\frac{\sqrt{k_x^2+k_y^2}}{R},\quad
  \phi=\operatorname{atan2}(k_y,k_x)
\end{equation}

\section{关节空间与操作空间正逆运动学}

根据常曲率假设，单关节中心线可近似为一段圆弧，曲率半径为
\begin{equation}
  \rho=\frac{L}{\theta}
\end{equation}
在弯曲平面内，末端相对近端的横向位移为 \(\rho(1-\cos\theta)\)，轴向高度为 \(\rho\sin\theta\)。再将该弯曲平面绕 \(z\) 轴旋转 \(\phi\)，可得到基坐标系下的末端位置：
\begin{equation}
  \left\{
  \begin{aligned}
    x &= \rho(1-\cos\theta)\cos\phi\\
    y &= \rho(1-\cos\theta)\sin\phi\\
    z &= \rho\sin\theta
  \end{aligned}
  \right.
\end{equation}
代入 \(\rho=L/\theta\)，得到关节空间到操作空间的正运动学：
\begin{equation}
  \left\{
  \begin{aligned}
    x &= \frac{L}{\theta}\cos\phi(1-\cos\theta)\\
    y &= \frac{L}{\theta}\sin\phi(1-\cos\theta)\\
    z &= \frac{L}{\theta}\sin\theta
  \end{aligned}
  \right.
\end{equation}
当 \(\theta\) 很小时，采用小角近似：
\begin{equation}
  x \approx \frac{L}{2}\theta\cos\phi,\quad
  y \approx \frac{L}{2}\theta\sin\phi,\quad
  z \approx L
\end{equation}

操作空间到关节空间的逆运动学为：已知末端位置 \(x,y,z\)，求 \(\theta,\phi\)。末端在 \(xy\) 平面内的径向距离为
\begin{equation}
  r_{xy}=\sqrt{x^2+y^2}
\end{equation}
由正运动学可得
\begin{equation}
  \frac{r_{xy}}{z}=\frac{\rho(1-\cos\theta)}{\rho\sin\theta}
  =\tan\frac{\theta}{2}
\end{equation}
因此
\begin{equation}
  \phi=\operatorname{atan2}(y,x),\quad
  \theta=2\operatorname{atan2}(r_{xy},z)
\end{equation}
得到 \(\theta,\phi\) 后，再代入 \(\Delta l_i=R\theta\cos(\phi-\alpha_i)\) 求三根绳长变化量，并由 \(q_i=\Delta l_i/r_w\) 求舵机角度，从而得到完整的操作空间到驱动空间逆解。

\section{末端姿态与齐次变换矩阵}

为描述末端姿态，将单关节弯曲运动分解为三个旋转：先绕 \(z\) 轴旋转 \(\phi\)，使弯曲平面与 \(xz\) 平面对齐；再绕局部 \(y\) 轴弯曲 \(\theta\)；最后绕局部 \(z\) 轴旋转 \(-\phi\)，回到原坐标方位。因此末端姿态矩阵为
\begin{equation}
  R_B^E = R_z(\phi)R_y(\theta)R_z(-\phi)
\end{equation}
其中
\begin{equation}
  R_z(\phi)=
  \begin{bmatrix}
    \cos\phi & -\sin\phi & 0\\
    \sin\phi & \cos\phi & 0\\
    0 & 0 & 1
  \end{bmatrix},
  \quad
  R_y(\theta)=
  \begin{bmatrix}
    \cos\theta & 0 & \sin\theta\\
    0 & 1 & 0\\
    -\sin\theta & 0 & \cos\theta
  \end{bmatrix}
\end{equation}
令 \(c_\phi=\cos\phi\)，\(s_\phi=\sin\phi\)，\(c_\theta=\cos\theta\)，\(s_\theta=\sin\theta\)，矩阵相乘后可得
\begin{equation}
  R_B^E =
  \begin{bmatrix}
    c_\phi^2c_\theta+s_\phi^2 & c_\phi s_\phi(c_\theta-1) & c_\phi s_\theta\\
    c_\phi s_\phi(c_\theta-1) & s_\phi^2c_\theta+c_\phi^2 & s_\phi s_\theta\\
    -c_\phi s_\theta & -s_\phi s_\theta & c_\theta
  \end{bmatrix}
\end{equation}

设末端位置向量为 \(p_B^E=[x,y,z]^T\)，则单关节末端相对于基坐标系的齐次变换矩阵为
\begin{equation}
  T_B^E =
  \begin{bmatrix}
    R_B^E & p_B^E\\
    0\quad 0\quad 0 & 1
  \end{bmatrix}
\end{equation}
展开为
\begin{equation}
  T_B^E =
  \begin{bmatrix}
    c_\phi^2c_\theta+s_\phi^2 & c_\phi s_\phi(c_\theta-1) & c_\phi s_\theta & x\\
    c_\phi s_\phi(c_\theta-1) & s_\phi^2c_\theta+c_\phi^2 & s_\phi s_\theta & y\\
    -c_\phi s_\theta & -s_\phi s_\theta & c_\theta & z\\
    0 & 0 & 0 & 1
  \end{bmatrix}
\end{equation}
其中 \(x,y,z\) 由关节空间到操作空间正运动学给出。该矩阵同时包含末端位置和姿态，可用于双关节串联时的坐标变换。

\section{图像误差到六舵机控制量的简化映射}

结合机构形式，本文采用“图像误差先滤波、再映射为舵机增量”的控制思路。记滤波后的误差为 \(\mathbf{e}_f(k)\)，则一阶滤波形式为
\begin{equation}
  \mathbf{e}_f(k)=\alpha\mathbf{e}(k)+(1-\alpha)\mathbf{e}_f(k-1),\quad 0<\alpha\leq 1
\end{equation}
记六个舵机的目标增量向量为 \(\Delta\mathbf{q}=[\Delta q_1,\Delta q_2,\Delta q_3,\Delta q_4,\Delta q_5,\Delta q_6]^T\)，可用待标定分配矩阵 \(\mathbf{M}\) 将二维图像误差映射为六维舵机增量：
\begin{equation}
  \Delta\mathbf{q}(k)=\operatorname{sat}\left(\mathbf{M}
  \begin{bmatrix}
    k_x & 0 \\
    0 & k_y
  \end{bmatrix}
  \mathbf{e}_f(k)\right)
\end{equation}
式中，\(k_x,k_y\) 为水平和竖直误差增益，\(\operatorname{sat}(\cdot)\) 表示对舵机目标增量、速度和累计位移进行限幅。分配矩阵 \(\mathbf{M}\) 由实体实验标定，用于确定图像误差方向、舵机旋转方向和拉索收放关系之间的对应关系。

\section{串口直控与安全约束}

舵机执行链路采用 PC 通过 USB/TTL 串口直接发送控制帧的方式实现。底层串口通信程序提供舵机扫描、基于时间的位置控制、基于速度的位置控制、阻尼模式、卸力释放模式和状态读取等功能。该接口能够直接下发目标角度并读取角度、电压、电流和功率状态，为闭环实验中的执行反馈记录提供基础。

为了避免软体机构在试验中出现过大拉索冲击，闭环实现至少需要满足三项约束：目标丢失时冻结或缓慢回零；对单次舵机增量和单位时间角度变化率设置上限；先在单关节模式下校正方向，再开展双关节协调实验。

\section{本章小结}

本章给出了图像误差到舵机控制量的基本链条：目标检测输出中心点，计算图像平面误差，经过滤波和限幅后映射为舵机增量，再由舵机牵引拉索产生关节弯曲。三绳单关节运动学公式和串口直控接口为控制实验提供了计算基础，误差方向、舵机分配矩阵和两关节职能分工可通过实体实验进一步标定。

\chapter{实验平台搭建与验证方案}

\section{实验目标与阶段安排}

实验平台验证目标包括三个方面：单目视觉侧是否能够完成图像采集、去畸变与目标误差提取；海参检测模型是否达到服务控制实验的可用水平；PC 直控舵机的执行链路是否能够支撑舵机联动控制。本章按视觉输入、检测跟踪、执行链路和闭环实验方案组织内容。

\begin{table}[H]
  \centering
  \caption{系统搭建状态}
  \label{tab:system-status}
  \begin{tabularx}{\textwidth}{>{\raggedright\arraybackslash}p{0.22\textwidth} >{\raggedright\arraybackslash}p{0.20\textwidth} X}
    \toprule
    模块 & 状态 & 在本实验平台中的作用 \\
    \midrule
    机械结构 & 已形成样机基础 & 完整设计为双关节六绳结构，实验先从第二关节验证展开 \\
    相机标定 & 已完成 & 51 幅有效棋盘格图像，平均重投影误差 0.0569 px \\
    目标检测与数据训练 & 已完成训练 & 单类别海参检测，60 幅图像、193 个标注框，为控制输入提供目标中心 \\
    误差提取 & 已具备接口 & 推理界面可输出目标中心相对图像中心的 \(dx,dy\) \\
    舵机驱动 & 已具备基础能力 & 串口扫描、角度下发、状态读取和卸力模式已实现 \\
    自动闭环 & 实验推进中 & 将图像误差映射到多舵机协调动作，并记录响应指标 \\
    \bottomrule
  \end{tabularx}
\end{table}

\section{视觉输入验证}

从相机标定结果来看，51 幅棋盘格图像全部被标定程序接受并用于参数求解，说明采集流程与角点检测较为稳定。结合去畸变预览图可以确认，标定结果已经进入视觉主流程，而不是停留在离线计算阶段。对于图像误差计算而言，这一步的意义在于统一图像几何基准，减少镜头畸变对目标中心位置的影响。

\section{检测与跟踪结果分析}

训练结果表明，YOLO26n 模型在小样本条件下具备较好的海参识别能力。最佳 mAP@0.5 为 0.993，最佳 mAP@0.5:0.95 为 0.742，说明模型在本验证集分布下能够较稳定地产生可用检测框。对视觉伺服任务而言，这些指标的意义主要体现在目标中心是否足够稳定，而不是追求更复杂数据集上的最终泛化结论。数据采集、标注、划分和训练共同构成本文视觉模块的实验工作，而不是单纯的前期准备。

进一步地，目标跟踪预览图能够给出目标中心与图像中心之间的偏差量，并显示实时帧率。这说明检测结果已经转化为控制器可使用的误差变量，视觉侧接口具备接入闭环控制的条件。

\section{执行链路与闭环实验设计}

在执行端，串口驱动脚本已经具备扫描舵机、读取状态和下发角度命令的功能，说明 PC 通过 USB/TTL 与舵机直接通信的方案可行。控制实验的关键环节，是将 \(e_x,e_y\) 的变化自动映射为舵机协调动作，并通过实体实验记录目标对准时间、稳态误差和动作平滑性等指标。

控制映射实验集中解决图像误差方向、舵机正反方向、拉索收放和软体关节弯曲之间的对应关系；实验量化记录对准时间、稳态像素误差、舵机角度变化和目标丢失恢复行为。这些指标用于评价视觉伺服控制是否能够稳定减小图像平面误差。

\section{本章小结}

本章基于实验图片、程序功能和训练结果，对系统平台进行了验证说明。机械平台、视觉输入和执行接口共同构成视觉伺服实验基础，闭环实验围绕响应时间、稳态误差和动作平滑性等指标展开。

\chapter{结论与展望}

\section{结论}

本文围绕海参采捕场景下双关节六绳驱动软体机器人的视觉伺服控制问题，完成了系统方案设计、视觉输入处理、目标检测训练、运动学建模和舵机执行链路搭建。与早期可伸缩吸取装置、ROS2 和 STM32 分层设想相比，本文根据实体样机和实验进度，将主线调整为“单目 USB 相机 + PC 通过 USB/TTL 直控舵机 + 绳驱软体机器人”的实验平台，并据此组织论文内容。

本文给出了双关节六绳驱动机构总体方案与样机基础，完成了单目相机标定与去畸变、海参检测数据集构建与模型训练、目标跟踪与图像误差提取、串口直控舵机接口验证等内容。结果表明，感知侧能够输出可供控制器使用的误差变量，执行侧具备直控基础，系统具备开展视觉伺服实验的条件。

在控制方法方面，本文建立了三绳单关节常曲率运动学模型，给出了驱动空间、关节空间和操作空间之间的映射关系，并提出了由图像误差生成舵机增量的控制链条。该模型适用于第一关节和第二关节各自的单关节分析，也为第一关节视觉对准、第二关节吸取范围扩展以及双关节协调控制提供了理论和程序实现基础。

\section{展望}

进一步工作围绕三个方向展开。第一，完成单关节和完整双关节条件下的方向校准实验，明确舵机与拉索在图像平面误差修正中的作用分工。第二，补充静态目标对准、动态跟踪和动作平滑性实验，记录响应时间、稳态误差和控制抖动等指标。第三，在 PC 直控链路稳定后，评估 ROS 化改造、模块化消息接口或更高层控制架构的必要性。

\chapter*{参考资料}

本文参考资料包括 Dewi 等关于模块化绳驱连续体机器人的论文、冯家琪关于水下绳驱连续型机械臂的学位论文、余文辉关于混合视觉伺服与模型预测控制的学位论文，以及曾洪杰关于视觉伺服与阻抗控制的学位论文。正式定稿时按学校模板统一整理为规范参考文献格式。

\end{document}
